% Part: first-order-logic
% Chapter: completeness
% Section: maximally-consistent-sets

% Definition of maximally consistent sets. Properties of maximally
% consistent sets required for the completeness proof are in
% maximally-consistent-sets.tex in the chapter on the proof system used.

\documentclass[../../../include/open-logic-section]{subfiles}

\begin{document}

\olfileid{fol}{com}{mcs}
\olsection{مجموعات \usetoken{P}{sentence} البالغة غاية الاتساق}

\begin{defn}[المجموعة المتسقة القصوى]
\ollabel{mcs}
تُسمى مجموعة~$\Gamma$ من~!!{sentence}s
\emph{متسقة قصوى} إذا وفقط إذا تحقق الشرطان الآتيان:
\begin{enumerate}
\item أن تكون~$\Gamma$ متسقة، و
\item أن يلزم من~$\Gamma \subsetneq \Gamma'$ كون~$\Gamma'$ غير متسقة.
\end{enumerate}
\end{defn}

\begin{explain}
وللحد السابق عبارة أخرى تكافئه: تُسمى مجموعة~$\Gamma$ من الجمل
\emph{متسقة قصوى} إذا وفقط إذا تحقق الآتي:
\begin{enumerate}
\item أن تكون~$\Gamma$ متسقة، و
\item أن يلزم من اتساق~$\Gamma \cup \{ !A \}$ كون~$!A \in \Gamma$.
\end{enumerate}
أي لا يسعنا أن نزيد مجموعة متسقة قصوى~!!{sentence}s ليست في~$\Gamma$
أصلاً، على مجموعة~$\Gamma$، إلا صارت المجموعة الأوسع الحاصلة غير متسقة.
\end{explain}

\begin{explain}
وإنما عظمت منزلة المجموعات المتسقة القصوى في برهان الاكتمال لأنا نضمن أن
كل مجموعة متسقة من~!!{sentence}s~$\Gamma$ داخلة في مجموعة متسقة قصوى
$\Gamma^*$؛ وأن هذه المجموعة القصوى تشتمل، لكل~!!{sentence}~$!A$، إما
على~$!A$ وإما على نفيها~$\lnot !A$. وهذا صادق، بوجه أخص، في
!!{sentence}s الذرية. فلنا، من مجموعة متسقة قصوى في لغة زيدت بما يلائم
من~!!{constant}s، أن ننشئ~!!a{structure} نحدّ فيها تأويل~!!{predicate}s
بما ينتمي إلى~$\Gamma^*$ من~!!{sentence}s الذرية. ثم يمكن إثبات أن هذه
!!{structure} تُصدق كل~!!{sentence}s في~$\Gamma^*$، فتُصدق من ثم ما في
$\Gamma$ أيضاً. ولا يتم برهان هذا إلا بأن يكون~$\lnot !A \in \Gamma^*$
إذا وفقط إذا كان~$!A \notin \Gamma^*$، وأن يكون
$(!A \lor !B) \in \Gamma^*$ إذا وفقط إذا كان~$!A \in \Gamma^*$ أو
$!B \in \Gamma^*$، وعلى هذا القياس.
\end{explain}

\begin{prop}
\ollabel{prop:mcs}
لتكن~$\Gamma$ متسقة قصوى؛ فعندئذ:
\begin{enumerate}
\item \ollabel{prop:mcs-prov-in} متى كان~$\Gamma \Proves !A$، كان
$!A \in \Gamma$.

\item \ollabel{prop:mcs-either-or} لكل~$!A$، يكون إما~$!A \in \Gamma$ أو
$\lnot !A \in \Gamma$.

\tagitem{prvAnd}{\ollabel{prop:mcs-and} يتحقق
$(!A \land !B) \in \Gamma$ إذا وفقط إذا تحقق كل من
$!A \in \Gamma$ و$!B \in \Gamma$.}{}

\tagitem{prvOr}{\ollabel{prop:mcs-or} يتحقق
$(!A \lor !B) \in \Gamma$ إذا وفقط إذا كان إما~$!A \in \Gamma$
وإما~$!B \in \Gamma$.}{}

\tagitem{prvIf}{\ollabel{prop:mcs-if} يتحقق
$(!A \lif !B) \in \Gamma$ إذا وفقط إذا كان إما~$!A \notin \Gamma$
وإما~$!B \in \Gamma$.}{}
\end{enumerate}
\end{prop}

\begin{proof}
نجعل~$\Gamma$، فيما يأتي كله، متسقة قصوى.
\begin{enumerate}
\item متى كان~$\Gamma \Proves !A$، كان~$!A \in \Gamma$.

افرض~$\Gamma \Proves !A$، ثم افرض، طلباً للتناقض، أن
$!A \notin \Gamma$. ولما كانت~$\Gamma$ متسقة قصوى، لزم أن تكون
$\Gamma \cup \{!A\}$ غير متسقة؛ ولذلك
$\Gamma \cup \{!A\} \Proves \lfalse$. وبحسب
\tagrefs{prfSC/{fol:seq:prv:prop:provability-contr},prfND/{fol:ntd:prv:prop:provability-contr}}،
تكون~$\Gamma$ غير متسقة، وهو خلاف ما فُرض من اتساقها. فيمتنع
$!A \notin \Gamma$، ويلزم~$!A \in \Gamma$.

\item لكل~$!A$، يكون إما~$!A \in \Gamma$ أو~$\lnot !A \in \Gamma$.

وافرض، ابتغاء التناقض، وجود~$!A$ يحقق~$!A \notin \Gamma$ و
$\lnot !A \notin \Gamma$. وبما أن~$\Gamma$ متسقة قصوى، تكون
$\Gamma \cup \{!A\}$ و$\Gamma \cup \{\lnot !A\}$ كلتاهما غير
متسقة؛ فيكون~$\Gamma \cup \{!A\} \Proves \lfalse$ و
$\Gamma \cup \{\lnot !A\} \Proves \lfalse$. وبحسب
\tagrefs{prfSC/{fol:seq:prv:prop:provability-exhaustive},prfND/{fol:ntd:prv:prop:provability-exhaustive}}،
تكون~$\Gamma$ غير متسقة، وذلك تناقض. فلا توجد~!!a{sentence}
$!A$ كذلك، ولكل~$!A$ يكون إما~$!A \in \Gamma$ أو
$\lnot !A \in \Gamma$.

\tagitem{defAnd}{}{%
\iftag{probAnd}{تمرين.}{%
يتحقق~$(!A \land !B) \in \Gamma$ إذا وفقط إذا تحقق كل من
$!A \in \Gamma$ و$!B \in \Gamma$:

فأما جهة اللزوم، فافرض~$(!A \land !B) \in \Gamma$. عندئذ
$\Gamma \Proves !A \land !B$. وبموجب
\tagrefs{prfSC/{fol:seq:prv:prop:provability-land-left},prfND/{fol:ntd:prv:prop:provability-land-left}}،
يكون~$\Gamma \Proves !A$ و$\Gamma \Proves !B$. ثم بحسب
\olref{prop:mcs-prov-in} يكون~$!A \in \Gamma$ و$!B \in \Gamma$، وهو
المقصود.

وأما جهة العكس، فليكن~$!A \in \Gamma$ و$!B \in \Gamma$. فيكون
$\Gamma \Proves !A$ و$\Gamma \Proves !B$. وبموجب
\tagrefs{prfSC/{fol:seq:prv:prop:provability-land-right},prfND/{fol:ntd:prv:prop:provability-land-right}}،
يكون~$\Gamma \Proves !A \land !B$. ثم من~\olref{prop:mcs-prov-in}
ينتج~$(!A \land !B) \in \Gamma$.}}

\tagitem{defOr}{}{%
\iftag{probOr}{تمرين.}{%
يتحقق~$(!A \lor !B) \in \Gamma$ إذا وفقط إذا كان إما
$!A \in \Gamma$ وإما~$!B \in \Gamma$.

ولإثبات نقيض جهة اللزوم، افرض~$!A \notin \Gamma$ و
$!B \notin \Gamma$. ونريد إثبات~$(!A \lor !B) \notin \Gamma$.
ولما كانت~$\Gamma$ متسقة قصوى، يكون
$\Gamma \cup \{!A\} \Proves \lfalse$ و
$\Gamma \cup \{!B\} \Proves \lfalse$. وبموجب
\tagrefs{prfSC/{fol:seq:prv:prop:provability-lor-left},prfND/{fol:ntd:prv:prop:provability-lor-left}}،
تكون~$\Gamma \cup \{(!A \lor !B)\}$ غير متسقة؛ فيلزم
$(!A \lor !B) \notin \Gamma$، وهو المراد.

وأما جهة العكس، فافرض~$!A \in \Gamma$ أو~$!B \in \Gamma$؛ فيكون
$\Gamma \Proves !A$ أو~$\Gamma \Proves !B$. وبحسب
\tagrefs{prfSC/{fol:seq:prv:prop:provability-lor-right},prfND/{fol:ntd:prv:prop:provability-lor-right}}،
يكون~$\Gamma \Proves !A \lor !B$. ثم من~\olref{prop:mcs-prov-in} ينتج
$(!A \lor !B) \in \Gamma$، وهو المطلوب.}}

\tagitem{defIf}{}{%
\iftag{probIf}{تمرين.}{%
يتحقق~$(!A \lif !B) \in \Gamma$ إذا وفقط إذا كان إما
$!A \notin \Gamma$ وإما~$!B \in \Gamma$:

فأما جهة اللزوم، فليكن~$(!A \lif !B) \in \Gamma$، ولنفرض ابتغاء
التناقض~$!A \in \Gamma$ و$!B \notin \Gamma$. ومن هذين الفرضين يكون
$\Gamma \Proves !A \lif !B$ و$\Gamma \Proves !A$. وبحسب
\tagrefs{prfSC/{fol:seq:prv:prop:provability-mp},prfND/{fol:ntd:prv:prop:provability-mp}}،
يكون~$\Gamma \Proves !B$. ولكن~\olref{prop:mcs-prov-in} تعطينا حينئذ
$!B \in \Gamma$، وذلك يناقض~$!B \notin \Gamma$.

وأما جهة العكس، فانظر أولاً في حالة~$!A \notin \Gamma$. بحسب
\olref{prop:mcs-either-or} يكون~$\lnot !A \in \Gamma$، فينتج
$\Gamma \Proves \lnot !A$. وبحسب
\tagrefs{prfSC/{fol:seq:prv:prop:provability-lif},prfND/{fol:ntd:prv:prop:provability-lif}}،
يكون~$\Gamma \Proves !A \lif !B$. وبتطبيق~\olref{prop:mcs-prov-in}
ثانية نحصل على~$(!A \lif !B) \in \Gamma$، وهو المطلوب.

ثم انظر في حالة~$!B \in \Gamma$. فعندئذ~$\Gamma \Proves !B$، وبحسب
\tagrefs{prfSC/{fol:seq:prv:prop:provability-lif},prfND/{fol:ntd:prv:prop:provability-lif}}،
يكون~$\Gamma \Proves !A \lif !B$. ومن~\olref{prop:mcs-prov-in} ينتج
$(!A \lif !B) \in \Gamma$.}}
\end{enumerate}
\end{proof}

\begin{prob}
أتمم برهان~\olref[fol][com][mcs]{prop:mcs}.
\end{prob}

\end{document}
